\documentclass{ltjsarticle}
\begin{document}
\section{定積分と不定積分}

関数$y=f(x)$のグラフと、$x=a, x=b\ (a<b)$の2直線と$x$軸で囲まれた領域の面積を$S$とするとき、$S$は以下のような\textbf{定積分}により求められます。
\[ S = \int_{a}^{b} f(x) dx \]
定積分の計算の中には、$f(x)$の\textbf{原始関数}$F(x)$が現れます。この原始関数を求める計算を効率的に表すために、\textbf{不定積分}がよく用いられます。
\[ \int f(x)dx = F(x) + C \]
ただし、$C$は積分定数です。

\section{さまざまな積分}
この積分という計算は他にもさまざまな種類に拡張されます。例えば、以下の2重積分では、$x=a,x=b, y=c,y=d$の4直線と$xy$平面、$z = f(x,y)$のグラフ(曲面)で囲まれた部分の体積$V$が求められます。
\[ V = \int_{c}^{d} \int_{a}^{b} f(x,y)dxdy \]

また、閉曲線$C$に沿って関数$w=f(z)$を連続的に足し合わせた値は、以下の周回積分により求められます。
\[ \oint_{C} f(z)dz \]
主に複素数関数論でこのような積分が現れます。
\end{document}

